\chapter{Building Structures}
\label{ch:building}

Everything in the \menu{Build} menu produces a new structure, either
replacing the active one or opening in a new tab. All of it runs through
ASE, so a working Python environment is a prerequisite throughout this
chapter.

\section{Supercells}

\menu{Edit}\menusep\menu{Unit Cell}\menusep\menu{Create Supercell} repeats
the active structure along its lattice vectors. Three spin boxes,
\ui{Repeat along a/b/c}, each 1--20, default $2 \times 1 \times 1$. The
operation is undoable and requires a defined cell along every repeated
direction.

\section{The surface slab wizard}
\label{sec:slab-wizard}

\menu{Build}\menusep\menu{Surface Slab} cleaves a slab from the active bulk
structure in three guided stages. Unlike a plain Miller-index dialog, each
stage shows you the consequence of your choice before you commit.

\subsection{Stage 1 --- surface orientation}

Set the Miller indices $(h\,k\,l)$ with three spin boxes ($-9$ to $9$), or
work the other way round: drag the red \textbf{u} and green \textbf{v}
handles on the interactive lattice canvas to any two lattice points, and
Calango computes the nearest integer $(h\,k\,l)$ for the plane they span.
The header shows the current vectors and resulting indices live, and the
info line reports $|u|$, $|v|$ and the angle between them.

Rotate the axonometric view by dragging; only non-collinear handle pairs
are accepted. \ui{Next} unlocks once a test cut succeeds --- $(0\,0\,0)$ is
rejected as not a plane.

\subsection{Stage 2 --- cut, thickness and terminations}

Calango builds a reference stack of eight bulk repeats and clusters the
atoms into atomic layers (0.10\,\AA{} tolerance). The cross-section canvas
draws each layer as a dashed line labelled with its $z$ coordinate, with
the current selection band highlighted between an orange \emph{top
termination} and a teal \emph{bottom termination}.

\begin{description}
  \item[Click a layer] Moves whichever boundary is nearer the click, so you
    can set either termination directly.
  \item[\ui{Layers in slab}] Sets the count numerically.
  \item[\ui{Thickness}] The same choice in \AA{}; it snaps to whole atomic
    layers above the bottom termination.
\end{description}

The info line reads back the selection, for example \emph{layers 1--8 of
16 (8 layers, 12.45\,\AA)}.

\subsection{Stage 3 --- vacuum and preview}

\ui{Top vacuum} and \ui{Bottom vacuum} (0--80\,\AA{}, default 10 each) pad
the cell along the surface normal. \ui{Symmetric vacuum (centered slab)},
checked by default, mirrors the top value to the bottom. A live 3D preview
shows the finished slab, with atom count, slab thickness and total cell
height. \ui{Finish} inserts it as a new tab labelled, for example,
\emph{(1 1 1) slab, 8 layers}.

\section{Liquid and gas interfaces}
\label{sec:liquid-interface}

\menu{Build}\menusep\menu{Liquid / Gas Interface} opens a fluid region on the
structure in the current tab and packs it with a liquid, a gas, a mixture, or
an ionic solution --- the solid/liquid and solid/gas cells that electro\-chem\-istry,
catalysis and adsorption studies start from. It consumes a structure rather
than generating one, so build or open the substrate first; a surface slab
from \S\ref{sec:slab-wizard} is the usual input. The result opens as a new
tab and the substrate's own tab is left untouched.

\subsection{Stage 1 --- geometry and region}

\begin{description}
  \item[\ui{Interface direction}] Which lattice vector the region is opened
    along; \emph{c} for a slab built by this application. The region is
    bounded by planes normal to the \emph{other} two lattice vectors, so it
    stays commensurate with the cell even when that cell is tilted.
  \item[\ui{Region thickness}] The gap between the top face of the structure
    and the bottom face of its own periodic image. The cell is grown --- or
    shrunk --- to make this \emph{exact}, whatever vacuum the input already
    carried. Asking for 20\,\AA{} of water and silently getting 20\,\AA{}
    plus the slab's existing 15\,\AA{} is the usual way to end up with an
    accidentally dilute interface.
  \item[\ui{Lateral supercell}] Repeats the substrate in the two directions
    that lie \emph{in} the interface plane, before the region is opened. A
    wider cell is usually what a solvation study needs: the fluid's
    correlation length is several \AA{}, and a $1\times1$ surface cell forces
    it to be periodic on a shorter length than that.
  \item[\ui{Surface clearance}] Closest approach between any fluid atom and
    any substrate atom (default 2.2\,\AA{}).
  \item[\ui{Move the substrate to the bottom of the cell}] On by default, so
    the fluid region is one contiguous block. Turn it off when the slab was
    positioned deliberately --- a symmetric slab centred in its cell, for
    instance.
\end{description}

The summary underneath reports the cell before and after, the fillable
volume, and how many water molecules that volume holds at ambient density.

\subsection{Stage 2 --- solvation and mixture}

\ui{Liquid / gas mixture} is a table of species and mole fractions. The
fractions are normalised, so $\{3\ \text{water},\ 1\ \text{ammonia}\}$ and
$\{0.75,\ 0.25\}$ request the same mixture. \ui{Amount} chooses between a
target mass density and an explicit molecule count.

\ui{Ionic species} inserts salts or bare ions by \emph{formula unit}. A salt
expands into its ions, so three units of (NH$_4$)$_2$SO$_4$ is six ammonium
and three sulfate and the cell stays neutral by construction. Bare ions are
offered too, for the cases where an unbalanced charge is deliberate; the net
charge is reported below the table, in orange when it is not zero.

Ions are placed \emph{before} the solvent and their mass counts against the
density target. Both matter: a packing that saturated could otherwise drop an
ion and silently change the stoichiometry, and a density target that ignored
the ions would give a brine at the density of water with salt added on top of
it.

\begin{caltip}
A gas at its 1\,bar density comes to a fraction of a molecule in a cell of
interface size, so picking a gas switches the page to an explicit molecule
count. That is the physically meaningful control: in a periodic cell the
molecule count \emph{is} the partial pressure.
\end{caltip}

\subsection{What the packer does, and what it does not}

Molecules are proposed at random positions with uniformly random
orientations and kept when they clear everything already placed. Two fluid
molecules are excluded on their centres, using a tabulated effective radius
per species, with an atom-level floor on top so that two molecules facing
each other cannot interlock their hydrogens; the substrate is checked atom by
atom, because that surface is fixed and a molecule started inside it cannot
relax out. The packing is seeded and therefore reproducible.

\begin{calwarn}
The result is a legitimate \emph{starting point} --- right composition, right
density, no overlaps --- and it is not an equilibrated liquid. A random
packing has no hydrogen-bond network and a potential energy far above the
equilibrium one. Quoting a diffusion coefficient, an interface energy or a
work of adhesion computed directly on this cell would be reporting a property
of the random number generator. Equilibrate with molecular dynamics
(\S\ref{sec:md}) first.
\end{calwarn}

Random sequential packing saturates below the equilibrium density of a
hydrogen-bonded liquid; around 90\% of the target is typical for water and is
perfectly usable. A larger shortfall is reported explicitly after the cell
opens, together with what to do about it.

\begin{calnote}
The species library holds only molecules whose geometry is fixed by symmetry
plus a bond length and, where needed, one angle: water, ammonia, HF, H$_2$S,
the common diatomic and small polyatomic gases, the noble gases, the alkali
and alkaline-earth cations, the halides, and the molecular ions
NH$_4^{+}$, H$_3$O$^{+}$, OH$^{-}$, NO$_3^{-}$, CO$_3^{2-}$, SO$_4^{2-}$ and
PO$_4^{3-}$. Species that need a real optimised geometry --- the alcohols,
anything with a torsional degree of freedom --- are deliberately absent
rather than approximated: a molecule with plausible-looking but wrong
internal coordinates produces a cell that nobody notices is wrong until the
first relaxation step tears it apart.
\end{calnote}

\section{Nanomaterials}

\menu{Build}\menusep\menu{Nanomaterial Builder} generates four families of
low-dimensional carbon and TMD structures. Pick the type at the top; the
form below changes to match.

\begin{description}
  \item[Graphene sheet] Lattice constant $a$ (default 2.46\,\AA), repeats
    $n_x \times n_y$ (default $4 \times 4$), vacuum (default 10\,\AA).
  \item[Graphene nanoribbon] Width and length in unit cells (default 6 each),
    \ui{Edge type} \ui{Zigzag} or \ui{Armchair}, and
    \ui{Hydrogen-terminate edges} (on by default).
  \item[Carbon nanotube] Chiral indices $n$ and $m$ (default 6, 6), length
    in unit cells, C--C bond length (default 1.42\,\AA), vacuum (default
    8\,\AA). $(0,0)$ is rejected.
  \item[TMD monolayer (MX$_2$)] An editable formula combo --- \texttt{MoS2},
    \texttt{MoSe2}, \texttt{WS2}, \texttt{WSe2}, \texttt{MoTe2},
    \texttt{NbSe2}, or anything else you type --- \ui{Phase} \texttt{2H} or
    \texttt{1T}, lattice constant (default 3.16\,\AA), X--X thickness
    (default 3.19\,\AA), repeats and vacuum.
\end{description}

\section{Metallic nanoparticles}
\label{sec:nanoparticle}

\menu{Build}\menusep\menu{Metallic Nanoparticle (Wulff)} builds clusters
two ways. Choose the element from the periodic table, then a mode.

\subsection{Wulff construction}

The equilibrium shape of a crystal is the one minimising total surface
energy at fixed volume; the Wulff construction produces it from the
relative surface energies of the exposed facets.

\begin{description}
  \item[Facet table] One row per facet: $h$, $k$, $l$ and
    $\gamma$ (relative surface energy). Defaults are $(111)\!=\!1.0$,
    $(100)\!=\!1.1$, $(110)\!=\!1.2$. \ui{Add Facet} and
    \ui{Remove Selected} edit the list.
  \item[\ui{Target atoms}] The requested size, default 200. The
    construction lands on the nearest achievable closed shape.
  \item[\ui{Size rounding}] \texttt{closest}, \texttt{above} or
    \texttt{below}.
  \item[\ui{Lattice}] \texttt{fcc}, \texttt{bcc} or \texttt{sc}.
\end{description}

Only the \emph{ratios} of the $\gamma$ values matter. Lowering
$\gamma(111)$ relative to $\gamma(100)$ grows the $(111)$ facets and drives
an fcc metal from a cube towards a truncated octahedron.

\subsection{Spherical cluster}

Carves every atom within a cutoff radius of the centre out of a bulk
lattice --- \texttt{fcc}, \texttt{bcc}, \texttt{sc} or \texttt{hcp}. Set
\ui{Radius} (default 10\,\AA). Useful when you want a specific diameter
rather than a thermodynamic shape.

Both modes centre the cluster with 6\,\AA{} of vacuum and clear periodicity.
\ui{Lattice constant} may be left at \ui{ASE default} to use the tabulated
value for the element.

\section{Special quasirandom structures}
\label{sec:sqs}

\menu{Build}\menusep\menu{Special Quasirandom Structure (SQS)} builds a
substitutional alloy whose short-range order approximates a truly random
solid solution --- the standard way to model a disordered alloy in a small
periodic cell.

\begin{description}
  \item[\ui{Supercell}] Repetitions $n_1 \times n_2 \times n_3$, default
    $2 \times 2 \times 2$.
  \item[\ui{Replace element}] Whose sites form the alloy sublattice.
  \item[\ui{Composition}] Target occupancies as symbol:fraction pairs, e.g.
    \texttt{Cu:0.75, Au:0.25}. Fractions are normalised and rounded to
    whole atoms by largest remainder.
  \item[\ui{First shell cutoff}, \ui{Second shell cutoff}] The pair shells
    whose Warren-Cowley parameters are driven toward zero. Defaults 3.2 and
    4.8\,\AA; set the second to \ui{off} to use one shell.
  \item[\ui{MC steps}, \ui{Random seed}] Anneal length and reproducibility.
\end{description}

If \texttt{icet} is importable, Calango uses its \texttt{generate\_sqs};
otherwise it falls back to a built-in simulated annealer that minimises
$\sum \alpha^2$ over the chosen shells using only NumPy and ASE neighbour
lists. The resulting tab is labelled with the backend that produced it, and
the internal backend also reports its residual objective --- a value near
zero means the decoration is essentially ideal.

\begin{caltip}
Verify the result with \menu{Analysis}\menusep\menu{Warren-Cowley analysis}
(Section~\ref{sec:warren-cowley}): a good SQS shows $\alpha \approx 0$ for
the shells you optimised.
\end{caltip}

\section{Normal modes and phonons}
\label{sec:phonon-builder}

\menu{Build}\menusep\menu{Normal Modes / Phonon Builder} sets up a
finite-displacement vibrational calculation. Calango detects whether the
structure is periodic and adapts:

\begin{description}
  \item[Periodic] Supercell finite displacements via \texttt{ase.phonons},
    giving dispersion and DOS.
  \item[Isolated molecule] $\Gamma$-point normal modes via
    \texttt{ase.vibrations}, giving frequencies and per-mode animation
    trajectories that Calango opens directly.
\end{description}

\subsection{Parameters}

\begin{description}
  \item[\ui{Mode}] \ui{Run vibrational calculation}, or
    \ui{Generate displaced structures only (new tab)} to export the
    displacement set for an external DFT code.
  \item[\ui{Supercell (a $\times$ b $\times$ c)}] Default $2\times2\times2$,
    periodic structures only.
  \item[\ui{Displacement $\delta$}] Default 0.01\,\AA.
  \item[\ui{Displaced structures}] A read-only count: $6N+1$ for $N$
    supercell atoms.
  \item[\ui{Calculator}] EMT, Lennard-Jones or MACE --- deliberately
    limited to calculators needing no external binary or pseudopotentials.
  \item[\ui{Band-path points}, \ui{DOS k-grid}] Sampling for the dispersion
    and the phonon density of states (defaults 100 and 20).
\end{description}

Like the calculator dialog, the generated ASE script is shown in an
editable pane and can be saved with \ui{Save Script\ldots}.

Outputs are $\Gamma$-point frequencies in the job console, plus
\file{phonon\_bands.csv} and \file{phonon\_dos.csv} in the job directory.

\begin{calnote}
Displacements are applied to every atom of the expanded supercell
($6N_\text{supercell}+1$ configurations), following the standard
finite-displacement recipe without symmetry reduction. For large
supercells this is the dominant cost; symmetry-reduced displacement sets
are a planned improvement.
\end{calnote}

\section{Structure perturbation and noise}
\label{sec:noise}

\menu{Build}\menusep\menu{Structure Perturbation / Noise} adds random
displacements --- for shaking a structure out of a symmetric saddle point,
or for generating training ensembles for machine-learning potentials.

\begin{description}
  \item[\ui{Distribution}] \ui{Gaussian (normal)} or \ui{Uniform}.
  \item[\ui{Amplitude}] Default 0.05\,\AA. For Gaussian this is $\sigma$
    per component; for uniform it is the half-width.
  \item[\ui{Random seed}] Default 42, for reproducibility.
  \item[\ui{Perturb atomic positions}] On by default.
  \item[\ui{Perturb unit cell vectors (random strain)}] Atoms follow the
    cell affinely, preserving fractional coordinates.
  \item[\ui{Frames}] 1 perturbs in place; more generates a trajectory.
  \item[\ui{Accumulation}] \ui{Independent} draws each frame from the
    original structure; \ui{Cumulative} performs a random walk from the
    previous frame.
\end{description}

A multi-frame run opens as a trajectory whose frame 0 is the unperturbed
original, is registered in the process manager, and is checkpointed to
\file{perturbed.extxyz} in the managed temporary directory
(Section~\ref{sec:calango-tmp}) --- so the ensemble can be reloaded later
and fed to the analysis or dataset tools.
